\chapterimage{chapter_head_1.pdf}
\chapter{Recursion \& Recursive Solutions}

\section{Theoretical Core Concepts}

\begin{definition}[Recursion]
Defining a problem in terms of a simpler version of itself.
\end{definition}

\begin{definition}[Base Case (Stopping Condition)]
The condition where the problem is small enough to be solved directly. It is required to prevent infinite recursion.
\end{definition}

\begin{definition}[Recursive Step (Reduction Step)]
Breaking the big problem into smaller subproblems by having the function call itself.
\end{definition}

\begin{itemize}
    \item \textbf{Stack Frames}: Each function call creates its own copy of all local variables, conceptually like stacking "Post-it notes". A return expression must be evaluated to a single value before control is passed back to the caller. 
    \item \textbf{Tail Recursion vs. Head Recursion}: Tail recursion is where the recursive call is the last operation. Head recursion is where the recursive call happens before any other operations.
\end{itemize}

\section{Algorithmic Tracing}

\begin{example}[Tracing factorial(4)]
Given the logic \texttt{return n * factorial(n - 1);} with base case \texttt{if (n <= 1) return 1;}
\begin{enumerate}
    \item \texttt{main()} calls \texttt{factorial(4)}
    \item Frame \texttt{f(4)}: evaluates \texttt{4 * f(3)}. Calls \texttt{f(3)}.
    \item Frame \texttt{f(3)}: evaluates \texttt{3 * f(2)}. Calls \texttt{f(2)}.
    \item Frame \texttt{f(2)}: evaluates \texttt{2 * f(1)}. Calls \texttt{f(1)}.
    \item Frame \texttt{f(1)}: hits base case \texttt{n <= 1}, returns \texttt{1}.
    \item Frame \texttt{f(2)}: receives \texttt{1}, evaluates \texttt{2 * 1 = 2}, returns \texttt{2}.
    \item Frame \texttt{f(3)}: receives \texttt{2}, evaluates \texttt{3 * 2 = 6}, returns \texttt{6}.
    \item Frame \texttt{f(4)}: receives \texttt{6}, evaluates \texttt{4 * 6 = 24}, returns \texttt{24} to \texttt{main()}.
\end{enumerate}
\end{example}

\section{Code Implementation}

\subsection{Towers of Hanoi}
\begin{lstlisting}[language=C]
void towerOfHanoi(int n, char source, char dest, char aux) {
    if (n == 1) {
        printf("Move disk 1 from pole %c to pole %c\n", source, dest);
        return;
    }
    // Move n-1 disks from source to aux, using dest as auxiliary
    towerOfHanoi(n - 1, source, aux, dest);
    // Move the nth disk from source to dest
    printf("Move disk %d from pole %c to pole %c\n", n, source, dest);
    // Move the n-1 disks from aux to dest, using source as auxiliary
    towerOfHanoi(n - 1, aux, dest, source);
}
\end{lstlisting}

\subsection{Permutations (Brute Force)}
\begin{lstlisting}[language=C]
void permutation_fill(Song *solution, int size, int pos, const Song *actual, int *is_used) {
    // Base Case: If the solution array is completely filled
    if (pos == size) {
        permutation_print(solution, size); 
        return;
    }    

    // Recursive Case: Fill the current slot and recursively fill remaining slots
    for (int x = 0; x < size; x++) {
        if (is_used[x] == 0) {
            solution[pos] = actual[x];
            is_used[x] = 1; // Mark element as used
            permutation_fill(solution, size, pos + 1, actual, is_used);
            is_used[x] = 0; // Backtrack (unmark)
        }
    }
}
\end{lstlisting}

\section{Skills \& Pitfalls}

\subsection{Common Mistakes}
\begin{remark}
\begin{itemize}
  \item \textbf{Infinite Recursion / Stack Overflow}: Forgetting a base case or passing parameters that do not progress toward the base case, eventually resulting in a Segmentation Fault.
  \item \textbf{Incorrect Order of Operations}: Putting operations after the recursive call when they should be before it (or vice versa), breaking intended logic.
\end{itemize}
\end{remark}

\subsection{Must-Know Skills}
\begin{itemize}
  \item \textbf{Brute Force Enumeration Framework}: Using recursion to systematically check all candidates for a solution (e.g., subsets, permutations). A \texttt{pos} index tracks depth, and a loop handles the choices at each index.
  \item \textbf{Tracking State}: Maintaining an \texttt{is\_used} array when generating permutations to ensure elements are not reused in the same candidate solution.
  \item \textbf{Understanding Overhead}: Recognizing that recursion tends to be less efficient than its iterative counterpart due to function call overhead. Some recursive algorithms (like basic Fibonacci) can do massive amounts of redundant work.
\end{itemize}
